% =====================================================================
% Paper C
% Closed-form exposure operator for coherent millimetre-wave
% beamforming on the human body
% Author: Robin Wydaeghe
% Target: IEEE Transactions on Wireless Communications
% =====================================================================

\documentclass[journal,twocolumn,10pt]{IEEEtran}

\usepackage[utf8]{inputenc}
\usepackage[T1]{fontenc}
\usepackage{lmodern}
\usepackage{microtype}

\usepackage{amsmath,amssymb,amsthm,mathtools}
\usepackage{bm}

\usepackage{graphicx}
\graphicspath{{figures/}{figures/paperC/}{../../theory/figures/}}
\usepackage{booktabs}
\usepackage{array}

\usepackage{cite}
\usepackage{xcolor}
\usepackage[colorlinks=true,linkcolor=blue,citecolor=blue,urlcolor=blue]{hyperref}
\usepackage[capitalize,nameinlink]{cleveref}

\theoremstyle{plain}
\newtheorem{theorem}{Theorem}
\newtheorem{proposition}[theorem]{Proposition}
\newtheorem{corollary}[theorem]{Corollary}
\theoremstyle{remark}
\newtheorem{remark}[theorem]{Remark}

\crefname{proposition}{Proposition}{Propositions}
\Crefname{proposition}{Proposition}{Propositions}
\crefname{corollary}{Corollary}{Corollaries}
\Crefname{corollary}{Corollary}{Corollaries}
\crefname{remark}{Remark}{Remarks}
\Crefname{remark}{Remark}{Remarks}
\crefname{theorem}{Theorem}{Theorems}
\Crefname{theorem}{Theorem}{Theorems}

\newcommand{\khat}{\hat{\bm{k}}}
\newcommand{\nhat}{\hat{\bm{n}}}
\newcommand{\rr}{\mathbf{r}}
\newcommand{\xx}{\mathbf{x}}
\newcommand{\hh}{\mathbf{h}}
\newcommand{\bvec}{\mathbf{b}}
\newcommand{\Sinc}{S_{\mathrm{inc}}}
\newcommand{\Sab}{S_{\mathrm{ab}}}
\newcommand{\GG}{\mathbf{G}}
\newcommand{\Gtil}{\tilde{\mathbf{G}}}
\newcommand{\Qmat}{\mathbf{Q}}
\newcommand{\Mmat}{\mathbf{M}}
\newcommand{\Jmat}{\mathbf{J}}
\newcommand{\Pmat}{\mathbf{P}}
\newcommand{\Phimat}{\mathbf{\Phi}}
\newcommand{\Psimat}{\bm{\Psi}}
\newcommand{\Psitil}{\tilde{\bm{\Psi}}}
\newcommand{\Fn}{\mathbf{F}_n}
\newcommand{\psin}{\bm{\psi}_n}
\newcommand{\diff}{\mathrm{d}}
\newcommand{\Complex}{\mathbb{C}}
\newcommand{\ntilde}{\tilde{n}}
\newcommand{\pospart}[1]{\left(#1\right)_{+}}
\newcommand{\Vis}{V}
\DeclareMathOperator{\ReLU}{ReLU}
\DeclareMathOperator{\RE}{Re}
\DeclareMathOperator{\IM}{Im}
\DeclareMathOperator{\tr}{tr}

\title{Closed-form exposure operator for coherent millimetre-wave
  beamforming on the human body}

\author{Robin~Wydaeghe%
  \thanks{R. Wydaeghe is with Ghent University and imec, Ghent, Belgium
  (e-mail: robin.wydaeghe@ugent.be).}}

\markboth{IEEE Transactions on Wireless Communications,
  Vol.~XX, No.~X, Month~2026}%
{Wydaeghe: Coherent exposure operator for mmWave beamforming}

\begin{document}

\maketitle

\begin{abstract}
A base station with $M$ antenna elements radiating with precoder
$\xx \in \Complex^M$ near a human body deposits whole-body absorbed
power as a Hermitian quadratic form $\xx^H \mathbf{S} \xx$. The
optimal exposure-constrained beamformer admits a closed-form solution
$\xx^\star \propto (\lambda \mathbf{S} + \nu \mathbf{I})^{-1} \hh^*$
for any positive-semidefinite $\mathbf{S}$. In every published
instance, $\mathbf{S}$ is calibrated by full-wave finite-difference
time-domain (FDTD) simulation per device per body configuration. At
millimetre-wave frequencies the cell count crosses $10^{12}$ and the
calibration becomes intractable. We derive the matrix in closed form
from the same propagation paths that give the user-equipment
channel $\hh$. Two bounded approximations, one on the refracted TM
polarisation direction with $\leq 4\,\%$ cross-term error and one on
the universal depth coupling with $\leq 0.5\,\%$ cross-term error,
collapse the coherent depth integral to a squared norm
$\Sab(\rr) = \|\Gtil(\rr)\xx\|^2$, where $\Gtil(\rr)$ is the
body-surface field channel. Integrating over the body gives the
exposure operator
$\Qmat = \int_\Sigma \Gtil(\rr)^H \Gtil(\rr)\,\diff A$. The operator
factors as $\Qmat = \Jmat^T \Mmat \Jmat$, where
$\Jmat \in \{0,1\}^{N \times M}$ is a combinatorial element-to-path
dispatch matrix and $\Mmat \in \Complex^{N \times N}$ is a Hermitian
path-pair Gram. The factorisation reduces the worst-case eigenvalue
to an $N$-dimensional path-space problem and admits broadband,
near-field, stochastic, low-rank, and multi-user extensions as small
perturbations. The closed-form Ying-2015 ECBF solution is retained,
with our analytical $\Qmat$ replacing the FDTD-calibrated SAR matrix.
A 64-element uniform rectangular array at 28\,GHz near the Thelonious
phantom shows a Pareto trade-off between received signal and absorbed
body power with a factor of 4 to 5 reduction in absorbed power at
50\,\% of MRT signal power. The exposure operator has effective rank
5 at 90\,\% trace energy regardless of multipath density.
\end{abstract}

\begin{IEEEkeywords}
Specific absorption rate, exposure-constrained beamforming, MIMO,
millimetre-wave, ray tracing, Fresnel transmission.
\end{IEEEkeywords}

\IEEEpeerreviewmaketitle

\section{Introduction}\label{sec:intro}
\IEEEPARstart{F}{ifth-generation} and emerging sixth-generation
wireless systems deploy multi-element antenna arrays at
millimetre-wave frequencies and operate them at short range from
human users. The 2020 ICNIRP guidelines~\cite{ICNIRP2020} change the
regulatory metric above 6\,GHz from specific absorption rate to
absorbed power density on the body surface, and tighten the local
limit at the same distances where mobile and fixed-wireless devices
operate. Exposure compliance has therefore moved from a post-hoc
verification step into the inner loop of array and waveform design.

The standard tool for exposure-aware beamforming is the SAR-matrix
formulation of Hochwald et al.~\cite{Hochwald2014}. The whole-body
absorbed power for an $M$-element array reduces to a Hermitian
quadratic form $\xx^H \mathbf{S} \xx$ in the precoder $\xx$.
Ying et al.~\cite{Ying2015} solved the exposure-constrained
signal-to-noise problem in closed form, showing that the optimal
precoder takes the regulariser-shifted form
$\xx^\star \propto (\lambda \mathbf{S} + \nu \mathbf{I})^{-1}\hh^*$.
Ebadi-Shahrivar et al.~\cite{Ebadi2019} reduced regulatory compliance
testing to a single eigenvalue $\lambda_{\max}(\mathbf{S})$. Variants
extend the approach to multi-user MIMO and robust precoding under
channel uncertainty~\cite{Ying2017}. This formulation is mature and
underpins every recent treatment of exposure-constrained MIMO in
the communications-theory literature.

The bottleneck is the matrix itself. In every
published instance, $\mathbf{S}$ is calibrated numerically by
full-wave finite-difference time-domain (FDTD) simulation on an
anatomical phantom at the operating frequency, with one calibration
per antenna mounting, per body posture, and per representative
array-body geometry. At sub-6\,GHz, the mesh requires of order
$10^9$ cells on a typical body and runs in tens of minutes per
configuration on a single GPU. At millimetre-wave, the cell count
climbs to $10^{12}$ per configuration and the calibration becomes
intractable inside design loops. Castellanos and
Heath~\cite{Castellanos2020} attack the cost at 28\,GHz by replacing
the matrix with per-element scalar Fresnel coefficients
$\beta \approx 0.70$. The reduction is calibrated against FDTD on a
fixed phantom-array geometry and discards angle-dependent phase
information that the coherent regime requires.

This paper derives the SAR-matrix analogue, the \emph{exposure
operator} $\Qmat$, in closed form from the same propagation paths
that ray tracing already produces. Two bounded approximations of
the near-tissue field collapse the body-surface absorbed power
density to a squared norm $\Sab(\rr) = \|\Gtil(\rr)\xx\|^2$, where
$\Gtil(\rr)$ is the body-surface field channel built
analytically from per-path Fresnel transmission. Integrating over
the body gives
$\Qmat = \int_\Sigma \Gtil(\rr)^H \Gtil(\rr)\,\diff A$
in $\Complex^{M \times M}$, with $\xx^H \Qmat \xx$ equal to the
whole-body absorbed power under precoder $\xx$. The two
approximations are exact on diagonal terms of $\Qmat$, so the
incoherent limit carries no error. The combined cross-term error
budget is $\leq 5\,\%$ for any precoder.

The operator factors in path space as
$\Qmat = \Jmat^T \Mmat \Jmat$, where
$\Jmat \in \{0,1\}^{N \times M}$ is the element-to-path dispatch
matrix and $\Mmat \in \Complex^{N \times N}$ is the Hermitian
path-pair Gram. The factorisation reduces the worst-case eigenvalue
computation to an $N$-dimensional path-space problem, and admits
broadband, near-field, stochastic, low-rank, and multi-user
extensions as local perturbations of $\Mmat$. The closed-form
exposure-constrained beamformer of Ying et al.~\cite{Ying2015}
is retained, with our analytical $\Qmat$ replacing the
FDTD-calibrated SAR matrix. A 64-element uniform rectangular array
at 28\,GHz near an anatomical phantom shows the construction in
practice. The Pareto trade-off between received signal and absorbed
body power gives a factor of 4 to 5 reduction in absorbed power at
50\,\% of maximum-ratio-transmission signal power, and the operator
has effective rank 5 at 90\,\% trace energy regardless of multipath
density.

The remainder of the paper is organised as follows.
\cref{sec:setup} fixes notation and assembles the
field, signal, and exposure quantities from a common ray-traced
channel. \cref{sec:coh-law} derives the squared-norm
absorption law. \cref{sec:path-space}
develops the path-space factorisation and its broadband,
near-field, stochastic, low-rank, and multi-user extensions.
\cref{sec:ecbf} treats the exposure-constrained beamforming
problem. \cref{sec:numerical} reports the 28\,GHz demonstration.


\section{System model}\label{sec:setup}

A base station has $M$ antenna elements and radiates with precoder
$\xx = [x_1, \ldots, x_M]^T \in \Complex^M$ such that
$\|\xx\|^2 = P$ is the total transmit power. The propagation
environment between the array and the body produces $N$ paths via
reflection, scattering, and diffraction. Path~$n$, for
$n = 1, \ldots, N$, is characterised by its originating element
$j(n) \in \{1, \ldots, M\}$, its direction of arrival $\khat_n$ at
the body surface, and its complex polarisation-amplitude vector
$\psin \in \Complex^3$ orthogonal to $\khat_n$. The polarisation-
amplitude vector encodes the path's amplitude at the body, the
polarisation state on arrival, and all propagation interactions
along the way (multipath fading, polarisation mixing through
reflectors, diffraction phase). Ray tracing engines such as Sionna
RT~\cite{SionnaRT} compute $\{j(n), \khat_n, \psin\}$ as their
natural output.

Three assumptions apply throughout. First, the body surface $\Sigma$
is locally flat on the scale of the wavelength, so Fresnel theory
applies pointwise. Second, the skin depth is much smaller than any
body dimension. At 28\,GHz, $\delta \sim 0.5\,$mm against body
curvature radii of $\sim 50\,$mm on the torso and $\sim 8\,$mm on
the smallest features (ears, fingers), so all transmitted power is
absorbed within a thin surface layer. Third, each path's last
scattering point sits at distance $d_n > 3\lambda$ from any
body-surface point of interest, so the local field at every
$\rr \in \Sigma$ is a superposition of locally plane waves. At
28\,GHz this requires $d_n > 3.2\,$cm, satisfied for any source
that is not in direct contact with the skin.

Three matrices organise the path-space machinery. The element-to-path
dispatch matrix is
\begin{equation}\label{eq:J-def}
  \Jmat \in \{0, 1\}^{N \times M},
  \qquad
  J_{nj} = \delta_{j(n), j}.
\end{equation}
Row $n$ has a single one in column $j(n)$. Column $j$ has a one in
each row corresponding to a path originating from element $j$. For
$\xx \in \Complex^M$, the lifted vector $\Jmat \xx \in \Complex^N$
broadcasts each element's weight onto all paths leaving that element.
For an observation point $\rr \in \mathbb{R}^3$, the phase diagonal
is
\begin{equation}\label{eq:Phi-def}
  \Phimat(\rr) = \mathrm{diag}\bigl(e^{-i k_0 \khat_1 \cdot \rr},
  \ldots, e^{-i k_0 \khat_N \cdot \rr}\bigr),
\end{equation}
a unitary diagonal that applies each path's geometric phase at
$\rr$. The polarisation matrix stacks the per-path
polarisation-amplitude vectors as columns,
\begin{equation}\label{eq:Psi-def}
  \Psimat = [\bm{\psi}_1, \ldots, \bm{\psi}_N] \in \Complex^{3 \times N}.
\end{equation}
$\Psimat$ depends on the propagation environment and the
transmit-antenna patterns. It does not depend on the precoder or the
observation point.

The free-space electric field at any point $\rr$ under precoder
$\xx$ is the sum of per-path contributions,
\begin{equation}\label{eq:E-decomp}
  \mathbf{E}(\rr)
  = \sum_n \psin\,x_{j(n)}\,e^{-i k_0 \khat_n \cdot \rr}
  = \GG(\rr)\,\xx,
\end{equation}
with field channel matrix
$\GG(\rr) = \Psimat\,\Phimat(\rr)\,\Jmat \in \Complex^{3 \times M}$.
The factorisation separates three concerns: the environment carried
by $\Psimat$, the position carried by $\Phimat$, and the dispatch
carried by $\Jmat$. The same factorisation recurs for the signal
and exposure channels. Define the path-amplitude carrier at $\rr$,
\begin{equation}\label{eq:carrier}
  \mathbf{c}(\rr, \xx) = \Phimat(\rr)\,\Jmat\,\xx \in \Complex^N.
\end{equation}
Entry $n$ of $\mathbf{c}$ is the precoder weight of path $n$'s
originating element times the path's geometric phase at $\rr$. The
carrier is the state variable from which every coherent observable
follows. The free-space field is $\mathbf{E}(\rr) = \Psimat\,
\mathbf{c}(\rr, \xx)$, the user-equipment received signal is
$y(\xx) = \bm\alpha^T \mathbf{c}(\rr_{\mathrm{UE}}, \xx) = \hh^T \xx$
with per-path antenna projections
$\alpha_n = \mathbf{C}_R(\khat_n)^H \psin$, equivalently
$\hh^* = \Jmat^T \bvec$ with
$\bvec = \Phimat(\rr_{\mathrm{UE}})^H \bm\alpha^*$. The
exposure-channel read-out is the subject of \cref{sec:coh-law}.


\section{Coherent absorption law}\label{sec:coh-law}

Two bounded approximations collapse the coherent depth integral to a
squared norm. At a surface point $\rr$ with normal $\nhat$, path $n$
has incidence cosine $\mu_n = \nhat \cdot (-\khat_n)$. For a
front-facing path ($\mu_n > 0$), the local TE and TM unit vectors are
$\hat{e}_{s,n} = \khat_n \times \nhat / |\khat_n \times \nhat|$ and
$\hat{e}_{p,n} = \hat{e}_{s,n} \times \khat_n$. The Fresnel
amplitude transmission coefficients are
\begin{equation}\label{eq:t-coef}
  t_{s,n} = \frac{2 \mu_n}{\mu_n + \xi_n},
  \quad
  t_{p,n} = \frac{2 \ntilde \mu_n}{\ntilde^2 \mu_n + \xi_n},
\end{equation}
with $\xi_n = \sqrt{\ntilde^2 - 1 + \mu_n^2}$ and
$\RE(\xi_n) > 0$. Writing $k_0 \xi_n = \beta_n - i\alpha_n^{(d)}$
with $\alpha_n^{(d)} > 0$, the depth factor inside tissue is
$e^{-i k_0 \xi_n z} = e^{-\alpha_n^{(d)} z} e^{-i \beta_n z}$. The
transmitted electric field of path $n$ at depth $z$ inside tissue is
\begin{equation}\label{eq:E-trans}
\begin{aligned}
  \mathbf{E}_n^{\mathrm{trans}}(\rr, z) ={}&
  \bigl(t_{s,n} \psi_{s,n} \hat{e}_{s,n} + t_{p,n} \psi_{p,n} \hat{e}_{p,n}\bigr) \\
  &\times x_{j(n)}\,e^{-i k_0 \khat_n \cdot \rr}\,e^{-i k_0 \xi_n z}\,
  \Vis(\rr,\khat_n)\,H(\mu_n),
\end{aligned}
\end{equation}
where $\psi_{s,n} = \hat{e}_{s,n} \cdot \psin$,
$\psi_{p,n} = \hat{e}_{p,n} \cdot \psin$, $H$ is the Heaviside step
(back-facing paths contribute zero), and $\Vis(\rr,\khat_n) \in
\{0,1\}$ is the binary self-shadowing visibility for path $n$ at
$\rr$, equal to one when $\rr$ is reached by an unobstructed ray
from $\khat_n$ and zero when shadowed by another body part. $\Vis$
is the ambient-occlusion factor of computer
graphics~\cite{Zhukov1998,Landis2002,AkenineMoller2018} and gates
self-shadowed contributions in the coherent sum. The total transmitted
field is the sum
$\mathbf{E}^{\mathrm{trans}} = \sum_n \mathbf{E}_n^{\mathrm{trans}}$,
and absorbed power density at $\rr$ is the ohmic loss integrated
over depth,
\begin{equation}\label{eq:Sab-int}
  \Sab(\rr) = \frac{\sigma}{2} \int_0^\infty
  |\mathbf{E}^{\mathrm{trans}}(\rr, z)|^2\,\diff z.
\end{equation}
Expanding the squared norm produces depth coupling factors
\begin{align}\label{eq:Lambda}
  \Lambda_{nn'}
  &= \int_0^\infty
    e^{-(\alpha_n^{(d)} + \alpha_{n'}^{(d)}) z}\,
    e^{i(\beta_{n'} - \beta_n) z}\,\diff z \notag\\
  &= \frac{1}{\alpha_n^{(d)} + \alpha_{n'}^{(d)} - i(\beta_{n'} - \beta_n)}.
\end{align}
The matrix $\Lambda$ is Hermitian positive-semidefinite. Diagonal
entries $\Lambda_{nn} = 1/(2\alpha_n^{(d)})$ recover the single-path
absorbed power. Off-diagonal entries carry the coherent cross-terms.

\subsection{Approximation 1: TM polarisation direction}

The exact transmitted TM field inside tissue points along the
refracted direction $\hat{e}'_{p,n}$, not the incident
$\hat{e}_{p,n}$. Snell's law gives
$\sin\theta_{t,n} = \sin\theta_n / |\ntilde|$, so the refracted
angle is $O(1/|\ntilde|)$ rather than $\theta_n$
(\cref{fig:incidence-plane-C}).

\begin{figure}[!t]
  \centering
  \includegraphics[width=\columnwidth]{incidence_plane_fig.pdf}
  \caption{Incidence-plane geometry for path $n$. The body surface
    $\Sigma$ has outward normal $\hat{\mathbf{n}}$. The incident ray
    arrives along $\hat{\mathbf{k}}_n$ at angle $\theta_n$, with
    polarisation amplitude vector $\boldsymbol{\psi}_n$ resolved on
    the TE basis $\hat{e}_{s,n}$ (out of plane) and TM basis
    $\hat{e}_{p,n}$ (in plane, large normal component
    $\sin\theta_n$). The refracted ray $\hat{\mathbf{k}}'_n$ enters
    at the much smaller angle $\theta_t = \arcsin(\sin\theta_n/|\ntilde|)$
    because of the high tissue refractive index, so the refracted
    TM direction $\hat{e}'_{p,n}$ is nearly tangential, with normal
    component $O(1/|\ntilde|)$.
    \cref{prop:approx1} replaces $\hat{e}'_{p,n}$ by $\hat{e}_{p,n}$
    in the cross-term sum. The discarded normal component is
    $\sin\theta_n(1 - 1/|\ntilde|)$ in amplitude, hence
    $O(1/|\ntilde|^2)$ in absorbed power.}
  \label{fig:incidence-plane-C}
\end{figure}

\begin{proposition}[TM polarisation direction]\label{prop:approx1}
Replace $\hat{e}'_{p,n}$ by $\hat{e}_{p,n}$ in
\cref{eq:E-trans}. Self-terms ($n = n'$) are unaffected because
$|\hat{e}_{p,n}|^2 = |\hat{e}'_{p,n}|^2 = 1$. The cross-term error is
$O(1/|\ntilde|^2) \leq 4\,\%$ for skin at 28\,GHz.
\end{proposition}

The mechanism is that the normal component of the electric field
inside tissue is $E_\perp \propto \sin\theta_t = O(1/|\ntilde|)$, so
its contribution to $|\mathbf{E}^{\mathrm{trans}}|^2$ is
$O(1/|\ntilde|^2)$. The bound holds uniformly in incidence angle.

\subsection{Approximation 2: universal depth coupling}

Define the normalised depth coupling
\begin{equation}\label{eq:Gamma-def}
  \Gamma_{nn'} \equiv \frac{\Lambda_{nn'}}{\sqrt{\Lambda_{nn} \Lambda_{n'n'}}}
  = \frac{2 \sqrt{\alpha_n^{(d)} \alpha_{n'}^{(d)}}}
  {\alpha_n^{(d)} + \alpha_{n'}^{(d)} - i(\beta_{n'} - \beta_n)}.
\end{equation}
By construction $\Gamma_{nn} = 1$. The cross-terms $\Gamma_{nn'}$
for $n \neq n'$ depend on the pair of incidence angles through
$\alpha_n^{(d)}$ and $\beta_n$.

\begin{proposition}[Universal depth coupling]\label{prop:approx2}
Set $\Gamma_{nn'} = 1$ for all path pairs. The cross-term error is
$|1 - \Gamma_{nn'}| \leq 0.44\,\%$ over
$(\theta_n, \theta_{n'}) \in [0^\circ, 85^\circ]^2$ at 28\,GHz on
skin, with mean $0.17\,\%$.
\end{proposition}

The mechanism is that the high refractive index $|\ntilde| \approx
4.83$ at 28\,GHz compresses all refraction angles to within
$\sim 10^\circ$ of normal, making $\alpha_n^{(d)}$ and $\beta_n$
nearly angle-independent.

\subsection{Reduction to a squared norm}

Define the per-path tissue-weighted polarisation read-out
\begin{equation}\label{eq:psi-tilde}
  \tilde{\bm{\psi}}_n(\rr) \equiv \sqrt{\frac{\sigma}{4\alpha_n^{(d)}}}\,
  \Fn(\rr)\,\psin \in \Complex^3,
\end{equation}
where $\Fn(\rr)$ is the per-path Fresnel transmission operator
that maps $\psin$ to $\Vis(\rr,\khat_n)\,(t_{s,n} \psi_{s,n} \hat{e}_{s,n} +
t_{p,n} \psi_{p,n} \hat{e}_{p,n})$ at incidence $\mu_n$ and surface
normal $\nhat(\rr)$, gating self-shadowed paths to zero. Stacking
$\Psitil(\rr) = [\tilde{\bm{\psi}}_1(\rr), \ldots,
\tilde{\bm{\psi}}_N(\rr)] \in \Complex^{3 \times N}$, the
body-surface field channel is
\begin{equation}\label{eq:Gtilde-decomp}
  \Gtil(\rr) = \Psitil(\rr)\,\Phimat(\rr)\,\Jmat \in \Complex^{3 \times M}.
\end{equation}
The structural parallel with $\GG(\rr) = \Psimat \Phimat \Jmat$
matters. Signal, field, and exposure channels share the dispatch
$\Jmat$ and the phase $\Phimat$. They differ only in the per-path
amplitude block: $\Psimat$ for the free-space field, $\bm\alpha^T$
for the user-equipment signal, $\Psitil$ for the body-surface
exposure. The signal and exposure branches are reciprocal in the
electromagnetic sense, with the signal branch sampling a single
point through the user-equipment antenna and the exposure branch
sampling the entire body surface through the per-path Fresnel
filter.

\begin{theorem}[Coherent absorption law]\label{thm:coh-law}
Under \cref{prop:approx1,prop:approx2},
\begin{equation}\label{eq:coh-law}
  \Sab(\rr) = \|\Gtil(\rr)\,\xx\|^2
  = \xx^H\,\Gtil(\rr)^H\,\Gtil(\rr)\,\xx,
\end{equation}
with $\Gtil(\rr) = \Psitil(\rr) \Phimat(\rr) \Jmat$ as in
\cref{eq:Gtilde-decomp}.
\end{theorem}

\begin{proof}
\Cref{prop:approx2} gives
$\Lambda_{nn'} = \sqrt{\Lambda_{nn}}\sqrt{\Lambda_{n'n'}}$, so the
expansion of $|\mathbf{E}^{\mathrm{trans}}|^2$ factors as
\begin{equation*}
\begin{aligned}
  \Sab(\rr)
  &= \frac{\sigma}{2} \sum_{n,n'} \Lambda_{nn'}\,
  [\mathbf{E}_n^{\mathrm{trans}}(\rr, 0)]^* \cdot
  \mathbf{E}_{n'}^{\mathrm{trans}}(\rr, 0) \\
  &= \biggl\|\sum_n \sqrt{\frac{\sigma}{4\alpha_n^{(d)}}}\,
  \mathbf{E}_n^{\mathrm{trans}}(\rr, 0)\biggr\|^2.
\end{aligned}
\end{equation*}
The sum inside the norm is the per-element contribution to
$\Gtil(\rr)\,\xx$ via \cref{eq:psi-tilde,eq:Gtilde-decomp}.
\end{proof}

\begin{corollary}[Single-wave limit]\label{cor:single-wave}
For a single path, \cref{eq:coh-law} reduces to the
polarisation-aware Fresnel surface law at a visible surface point,
and under the additional pseudo-Brewster approximation, to the
geometric law $\Sab(\rr) = \Sinc T_0 \Vis(\rr,\khat) \pospart{\mu}$.
\end{corollary}

\begin{corollary}[Incoherent limit]\label{cor:incoherent}
When the $N$ paths have mutually random relative phases, cross-terms
average to zero and the expected absorbed power density reduces to
$\langle\Sab(\rr)\rangle = Z_0^{-1}\sum_n |x_{j(n)}|^2 (T_{s,n}
|\psi_{s,n}|^2 + T_{p,n} |\psi_{p,n}|^2) \Vis(\rr,\khat_n)
\pospart{\mu_n}$, the incoherent multi-source surface dosimetry
formula.
\end{corollary}

The combined error budget on $\Sab$ is the sum of the two
cross-term contributions, $\leq 5\,\%$ on $\xx^H \Qmat \xx$ for any
precoder. Self-terms are exact, so the incoherent limit carries no
approximation error.


\section{Path-space factorisation}\label{sec:path-space}

Integrating \cref{eq:coh-law} over the body surface gives a
Hermitian quadratic form in the precoder. Define the exposure
operator
\begin{equation}\label{eq:Q-def}
  \Qmat = \int_\Sigma \Gtil(\rr)^H\,\Gtil(\rr)\,\diff A
  \in \Complex^{M \times M},
\end{equation}
so that $P_{\mathrm{abs}} = \xx^H\,\Qmat\,\xx$ is the whole-body
absorbed power under precoder $\xx$. $\Qmat$ is Hermitian
positive-semidefinite by construction, with rank at most
$\min(M, 3 M_\Sigma)$ where $M_\Sigma$ is the number of integration
triangles on the body mesh. The eigendecomposition
$\Qmat = \sum_k \lambda_k \mathbf{v}_k \mathbf{v}_k^H$ with
$\lambda_1 \geq \cdots \geq \lambda_M \geq 0$ are the exposure
modes. The leading eigenvector $\mathbf{v}_1$ is the precoder
direction that maximises body absorption per unit transmit power,
and $\lambda_1$ is that maximum. The trace
$\tr(\Qmat) = \sum_k \lambda_k$ equals the expected absorbed power
under isotropic random precoding, scaled by the transmit power.

The Hochwald-Ying SAR matrix $\mathbf{S}$ is structurally identical
to $\Qmat$. They differ in three ways. First, $\mathbf{S}$ is
calibrated by FDTD on a specific phantom-array configuration, while
$\Qmat$ is computed in closed form from propagation paths. Second,
$\mathbf{S}$ is defined as a volume integral over a fixed tissue
volume $V$, while $\Qmat$ is the surface integral over the entire
body, giving whole-body absorbed power directly. Third, the entries
of $\mathbf{S}$ have no closed-form physical content, while the
entries of $\Qmat$ factor through the Fresnel coefficients and the
path geometry as shown below. The leading-eigenvalue identity
$\lambda_{\max}(\mathbf{S}) = \max_x \mathrm{SAR}(\xx) / \|\xx\|^2$
of \cite{Ebadi2019} carries over directly: $\lambda_{\max}(\Qmat)$
is the worst-case whole-body absorbed power per unit transmit
power. The Castellanos per-element scalar $\beta \approx 0.70$
\cite{Castellanos2020} is recovered as a special case of
\cref{eq:psi-tilde}: at normal incidence, each path contributes
$\sqrt{T_0/(2\alpha_0^{(d)} Z_0)} = \beta_{\mathrm{Castellanos}}$.
The per-path operator $\Fn(\rr)$ retains the angle-dependent phase
that the scalar reduction discards.

The structural identity
$\Gtil(\rr) = \Psitil(\rr) \Phimat(\rr) \Jmat$ makes the surface
integral defining $\Qmat$ separate cleanly. The $N \times N$
Hermitian positive-semidefinite path-pair Gram is
\begin{equation}\label{eq:M-def}
  \Mmat = \int_\Sigma \Phimat(\rr)^H\,\Psitil(\rr)^H\,
  \Psitil(\rr)\,\Phimat(\rr)\,\diff A.
\end{equation}
Entry $(n, n')$ is
\begin{equation}\label{eq:M-entry}
\begin{aligned}
  M_{nn'} = {}& \sqrt{\frac{\sigma^2}{16\,\alpha_n^{(d)}\alpha_{n'}^{(d)}}}\,
  \int_\Sigma \psi_n^H\,\Fn(\rr)^H\,\mathbf{F}_{n'}(\rr)\,\psi_{n'} \\
  &\times\,e^{-i k_0 (\khat_{n'} - \khat_n) \cdot \rr}\,\diff A.
\end{aligned}
\end{equation}
The phase kernel inside the integral is $e^{-i \mathbf{q}_{nn'}
\cdot \rr}$ with $\mathbf{q}_{nn'} = k_0(\khat_{n'} - \khat_n)$.
Off-diagonal entries of $\Mmat$ are body-surface Fourier components
at the difference wave vector of each path pair, weighted by the
inner product of the Fresnel-filtered polarisations. Diagonal
entries reduce to the single-path absorbed-power surface integral.

\begin{theorem}[Path-space factorisation]\label{thm:Q-factor}
The exposure operator factors as
\begin{equation}\label{eq:Q-factor}
  \Qmat = \Jmat^T\,\Mmat\,\Jmat.
\end{equation}
\end{theorem}

\begin{proof}
From $\Gtil(\rr) = \Psitil(\rr)\Phimat(\rr)\Jmat$ and linearity of
the integral,
\begin{equation*}
  \Qmat
  = \Jmat^T \!\!\int_\Sigma \Phimat(\rr)^H\,\Psitil(\rr)^H\,
  \Psitil(\rr)\,\Phimat(\rr)\,\diff A\,\Jmat
  = \Jmat^T\,\Mmat\,\Jmat. \qedhere
\end{equation*}
\end{proof}

The factorisation isolates two distinct types of object. The
combinatorial dispatch $\Jmat$ records which transmit element drives
which path. It depends on the array hardware and on the ray-tracing
output but has no continuous structure. The path-pair Gram $\Mmat$
is the entire continuous content of the body's response. It depends
on the path geometry, the body shape, and the tissue dielectric
properties, but not on the array hardware.

\subsection{Worst-case eigenvalue lift}

Define the element-block projector $\Pmat = \Jmat \Jmat^T \in
\{0,1\}^{N \times N}$. With paths grouped by originating element,
$\Pmat$ is block-diagonal. Each block is an all-ones matrix of size
$N_j \times N_j$, where $N_j$ is the number of paths originating
from element $j$.

\begin{corollary}[Worst-case eigenvalue lift]\label{cor:lambda-lift}
The nonzero eigenvalues of $\Qmat = \Jmat^T \Mmat \Jmat$ coincide
with those of $\Mmat \Pmat$. In particular,
\begin{equation}\label{eq:lambda-lift}
  \lambda_{\max}(\Qmat) = \lambda_{\max}(\Mmat\,\Pmat).
\end{equation}
A worst-case precoder $\xx_\star$ lifts from the path-space
eigenvector $\mathbf{u}_\star \in \Complex^N$ of $\Mmat \Pmat$ via
$\xx_\star \propto \Jmat^T \mathbf{u}_\star$.
\end{corollary}

\begin{proof}
For matrices $A \in \Complex^{p \times q}$ and $B \in \Complex^{q
\times p}$, the nonzero spectra of $AB$ and $BA$ coincide. Apply
with $A = \Jmat^T$ and $B = \Mmat \Jmat$: $\Jmat^T (\Mmat \Jmat) =
\Qmat$ and $(\Mmat \Jmat) \Jmat^T = \Mmat \Pmat$.
\end{proof}

The element-tied subspace $\mathcal{E} = \mathrm{range}(\Jmat)
\subset \Complex^N$ is the geometric locus of reachable absorption
modes. Path-space directions orthogonal to $\mathcal{E}$ cannot be
excited by any precoder. The dimension of $\mathcal{E}$ is at most
$M$, so the spectral problem on $\Mmat \Pmat$ has at most $M$
nonzero eigenvalues regardless of $N$. When $N \ll M$ the
path-space problem is computationally cheaper. When $N \gg M$, the
projector $\Pmat$ enforces the rank limit and the relevant subspace
remains $M$-dimensional.

\subsection{Extensions}

The factorisation makes four extensions immediate.

\emph{Broadband.} For a precoder $\xx(\omega) \in \Complex^M$ with
finite transmit energy $\int |\xx(\omega)|^2\,\diff\omega/(2\pi) \leq
P$, time-averaged absorbed power is
\begin{equation*}
  \bar{P}_{\mathrm{abs}} = \int \xx(\omega)^H \Qmat(\omega)
  \xx(\omega)\,\frac{\diff\omega}{2\pi},
  \;\;
  \Qmat(\omega) = \Jmat^T \Mmat(\omega) \Jmat.
\end{equation*}
The factorisation holds at each frequency. The broadband
exposure-constrained problem decouples pointwise in $\omega$ with
shared dual multipliers across the integrated constraints.

\emph{Near-field.} When the last-scatter distance violates
$d_n > 3\lambda$, replace the unitary $\Phimat(\rr)$ with a confocal
phase-and-amplitude kernel $\Phimat^{\mathrm{NF}}(\rr) =
\mathrm{diag}(\varphi_n(\rr))$ where $\varphi_n(\rr) =
(R_n^{(0)}/R_n(\rr))\,e^{-i k_0 (R_n(\rr) - R_n^{(0)})}$ and
$R_n(\rr)$ is the distance from the path's last-scatter point to
$\rr$. The carrier factorisation
$\mathbf{c}^{\mathrm{NF}}(\rr, \xx) = \Phimat^{\mathrm{NF}}(\rr)
\Jmat \xx$ retains the same shape, and the worst-case eigenvalue
lift proceeds verbatim.

\emph{Stochastic channels.} For a random $\Mmat$ induced by
stochastic path parameters with $\bar{\Mmat} = \mathbb{E}[\Mmat]$,
linearity of the integral gives $\mathbb{E}[\Qmat] = \Jmat^T
\bar{\Mmat} \Jmat$. The expected-value exposure-constrained problem
inherits the closed-form solution with $\Mmat$ replaced by
$\bar{\Mmat}$.

\emph{Low-rank approximation.} For rank-$r$ truncation $\Qmat^{(r)}$,
the worst-case error on $\xx^H \Qmat \xx$ is $\lambda_{r+1} P$. The
top-$r$ eigenvalues live on $\mathcal{E} \subset \Complex^N$ and can
be computed by Lanczos on $\Mmat \Pmat$ at cost $O(r \|\Mmat\|_0)$
per iteration, where $\|\Mmat\|_0$ counts non-negligible entries of
$\Mmat$. For typical multipath at mmWave, $r \sim 10$ accounts for
the dominant absorption modes regardless of $M$.


\section{Exposure-constrained beamforming}\label{sec:ecbf}

The optimal precoder under a body-absorption constraint admits a
closed-form solution that lifts to path space. The standard problem
maximises received signal power subject to a body-absorption budget,
\begin{equation}\label{eq:ECBF}
  \max_\xx |\hh^T \xx|^2
  \;\;\text{s.t.}\;\;
  \xx^H\,\Qmat\,\xx \leq P_{\mathrm{abs}}^{\max},
  \;\;\|\xx\|^2 \leq P.
\end{equation}
Both constraints are convex. The objective is the squared modulus of
a linear functional of $\xx$. The problem is a non-convex QCQP with
two PSD constraints, for which the S-procedure guarantees that the
KKT point is a global maximum~\cite{BoydVandenberghe2004,Ying2015}.

\begin{theorem}[Optimal precoder]\label{thm:ecbf}
The maximum of \cref{eq:ECBF} is attained at
\begin{equation}\label{eq:x-opt}
  \xx^\star = \sqrt{P}\,
  \frac{(\lambda \Qmat + \nu \mathbf{I})^{-1} \hh^*}
  {\|(\lambda \Qmat + \nu \mathbf{I})^{-1} \hh^*\|},
\end{equation}
with Lagrange multipliers $\lambda \geq 0$ and $\nu \geq 0$
determined by complementary slackness.
\end{theorem}

This is the form derived by Ying et al.~\cite{Ying2015} for the
FDTD-calibrated SAR matrix, with our analytical $\Qmat$
substituted. The matrix $\lambda \Qmat + \nu \mathbf{I}$ is positive
definite whenever $\nu > 0$, which is the generic case (the power
constraint is generically active).

\begin{theorem}[Path-space push-through]\label{thm:push-through}
The optimal precoder lifts to path space as
\begin{equation}\label{eq:x-path}
  \xx^\star \propto \Jmat^T\,
  (\nu \mathbf{I}_N + \lambda \Mmat \Pmat)^{-1}\,\bvec,
\end{equation}
with $\bvec = \Phimat(\rr_{\mathrm{UE}})^H \bm\alpha^*$ the
path-space signal vector.
\end{theorem}

\begin{proof}
Apply the push-through identity $(\mathbf{I}_M + AB)^{-1} A = A
(\mathbf{I}_N + BA)^{-1}$ with $A = \Jmat^T$ and $B = (\lambda/\nu)
\Mmat \Jmat$. Then $AB = (\lambda/\nu) \Qmat$ and $BA = (\lambda/\nu)
\Mmat \Pmat$, so
\begin{equation*}
  (\mathbf{I}_M + (\lambda/\nu) \Qmat)^{-1} \Jmat^T
  = \Jmat^T (\mathbf{I}_N + (\lambda/\nu) \Mmat \Pmat)^{-1}.
\end{equation*}
Multiplying by $1/\nu$ and applying to $\bvec$ gives the result.
\end{proof}

The element-space inverse in \cref{eq:x-opt} is $M \times M$. The
path-space inverse in \cref{eq:x-path} is $N \times N$. When
$N < M$, the path-space form is computationally cheaper. When
$N > M$, the rank of $\Mmat \Pmat$ is at most $M$ and a rank-$r$
Lanczos iteration suffices.

The exposure constraint slack ($\lambda \to 0$) recovers
$\xx^\star \propto \hh^*$, standard MRT. The exposure constraint
tight ($\lambda \to \infty$ with $\nu$ finite) gives
$\xx^\star \approx \Qmat^{-1} \hh^*$ when $\Qmat$ is full rank, or
the pseudoinverse otherwise. The precoder steers into the absorption
null space of the body. The signal-per-absorbed-power generalised
eigenvalue problem ignoring the power constraint has the closed
form $\xx^\star \propto \Qmat^{-1} \hh^*$ with value
$\hh^T \Qmat^{-1} \hh^*$. This traces the Pareto-optimal trade-off
boundary in the limit $P \to \infty$.

Under MRT, the precoder is $\xx_{\mathrm{MRT}} = \sqrt{P}\,\hh^* /
\|\hh\|$ and the signal power is $|y|^2 = P \|\hh\|^2$. The body
absorbed power under MRT is
\begin{equation}\label{eq:Pabs-MRT}
  P_{\mathrm{abs}}^{\mathrm{MRT}}
  = \frac{P}{\|\hh\|^2}\,\hh^T\,\Qmat\,\hh^*.
\end{equation}
The exposure-signal alignment is
\begin{equation}\label{eq:rho-def}
  \rho = \frac{\hh^T\,\Qmat\,\hh^*}{\|\hh\|^2\,\lambda_{\max}(\Qmat)}
  \in [0, 1],
\end{equation}
with $P_{\mathrm{abs}}^{\mathrm{MRT}} = P\,\rho\,\lambda_{\max}(\Qmat)$.
The bound $\rho \in [0, 1]$ is the Rayleigh quotient applied to
$\Qmat$ at $\hh^*$. The value $\rho = 1$ means MRT is the worst-case
absorption direction, while $\rho \ll 1$ means the dominant signal
paths are misaligned with the body's absorption modes and the body
is geometrically protected from MRT. In path space, using
$\hh^* = \Jmat^T \bvec$,
\begin{equation}\label{eq:rho-path}
  \rho = \frac{\bvec^H\,\Pmat\,\Mmat\,\Pmat\,\bvec}
  {(\bvec^H \Pmat \bvec)\,\lambda_{\max}(\Mmat\,\Pmat)},
\end{equation}
which makes the geometric content explicit. Three mechanisms
decouple body hot spot from UE focus. First, the signal sums scalar
projections $h_j^* h_j \in \mathbb{R}_+$ while the body field sums
vector contributions in $\Complex^3$. Even when all
magnitudes match, the vectors may partially cancel when transmitted
polarisations differ at $\rr$. Second, the UE channel projects each
path through a scalar antenna inner product, while the exposure
channel applies a $3 \times 3$ rank-2 Fresnel operator that depends
on the surface normal. A grazing path may contribute strongly to
the signal but weakly to absorption because both Fresnel coefficients
vanish at $\mu = 0$. Third, the signal branch is evaluated at a
single point where carrier phases are aligned by MRT design, while
the exposure branch is evaluated at every $\rr \in \Sigma$ where
each path acquires an additional phase $e^{-i k_0 \khat_n \cdot
(\rr - \rr_{\mathrm{UE}})}$. Coherent alignment that holds at the
UE need not hold at the body.


\section{Multi-user MIMO}\label{sec:mu}

The single-user form extends to $K$ simultaneous users without
new approximations. In a multi-user MIMO downlink, the base station
transmits to $K$ users with data symbols $s_k$ (zero-mean, unit
variance, mutually uncorrelated) and precoding vectors $\mathbf{w}_k
\in \Complex^M$. The total transmitted signal is $\xx_{\mathrm{total}}
= \sum_{k=1}^K \mathbf{w}_k s_k$ with total power $\sum_k
\|\mathbf{w}_k\|^2 \leq P$. The propagation environment produces a
single set of paths shared across all users. The body of person~$u$
in the scene has surface $\Sigma^{(u)}$ and per-person exposure
channel $\Gtil^{(u)}(\rr) = \Psitil^{(u)}(\rr) \Phimat(\rr) \Jmat$
for $\rr \in \Sigma^{(u)}$. The path geometry $(\Phimat, \Jmat)$ is
shared across all bodies. The per-person Fresnel filter
$\Psitil^{(u)}$ depends on each body's normals and tissue properties.
The per-person exposure operator is
\begin{equation}\label{eq:Q-u}
  \Qmat^{(u)} = \int_{\Sigma^{(u)}} \Gtil^{(u)}(\rr)^H\,
  \Gtil^{(u)}(\rr)\,\diff A
  = \Jmat^T\,\Mmat^{(u)}\,\Jmat,
\end{equation}
with $\Mmat^{(u)}$ the per-person path-pair Gram. Each $\Qmat^{(u)}$
is Hermitian PSD on $\Complex^{M \times M}$.

The absorbed power density on body~$u$ is
$S_{\mathrm{ab}}^{(u)}(\rr) = \|\Gtil^{(u)}(\rr)
\xx_{\mathrm{total}}\|^2$. Taking expectations over the data
symbols, the cross-stream terms vanish:
\begin{equation}\label{eq:Sab-MU}
  \mathbb{E}[S_{\mathrm{ab}}^{(u)}(\rr)]
  = \sum_{k=1}^K \|\Gtil^{(u)}(\rr)\,\mathbf{w}_k\|^2.
\end{equation}
The time-averaged total absorbed power on body~$u$ is
$P_{\mathrm{abs}}^{(u)} = \sum_{k=1}^K \mathbf{w}_k^H\,\Qmat^{(u)}\,
\mathbf{w}_k$. The standard MU-MIMO objective maximises sum-rate.
With per-person exposure constraints,
\begin{equation}\label{eq:MU-ECBF}
\begin{aligned}
  \max_{\{\mathbf{w}_k\}} {}&\sum_{k=1}^K \log_2\!\biggl(1 +
  \frac{|\hh_k^T \mathbf{w}_k|^2}{\sum_{k' \neq k}
  |\hh_k^T \mathbf{w}_{k'}|^2 + \sigma_n^2}\biggr) \\
  \text{s.t. } & P_{\mathrm{abs}}^{(u)} \leq P_{\mathrm{abs}}^{\max},
  \quad u = 1, \ldots, U, \\
  & \textstyle\sum_k \|\mathbf{w}_k\|^2 \leq P.
\end{aligned}
\end{equation}
The exposure constraints are quadratic in the $\{\mathbf{w}_k\}$
with PSD matrices $\Qmat^{(u)}$, and integrate directly into
existing weighted-MMSE solvers~\cite{Christensen2008,Shi2011}. No
additional approximation is required.


\section{Numerical demonstration}\label{sec:numerical}

The construction above is exercised on a representative scenario at
28\,GHz. An $8 \times 8 = 64$-element uniform rectangular array with
half-wavelength spacing is centred at $\mathbf{r}_{\mathrm{BS}} =
(0, 0, 1.5)\,$m, the Thelonious anatomical phantom is centred at
the origin, and two user-equipment positions are placed at
$\rr_{\mathrm{UE}}^{(a)} = (0, 5, 1)\,$m and
$\rr_{\mathrm{UE}}^{(b)} = (0, -5, 1)\,$m. Multipath is generated by
clustered scattering with 40 paths per transmit element distributed
on a $15^\circ$ cone around the line-of-sight direction, giving $N
= 2560$ total paths. Per-element radiation patterns are isotropic.
Skin tissue at 28\,GHz is taken from the IT'IS database (v5.0):
$\varepsilon_r = 16.55$, $\sigma = 25.8\,$S/m,
$\ntilde = 4.47 - 1.85i$.

\subsection{Pareto trade-off}

For each UE position the Lagrange multiplier $\lambda$ is swept over
50 values logarithmically spaced from $10^{-3}$ to $10^3$ at fixed
$\nu = 1$, and received signal power $|\hh^T \xx|^2/P$ is plotted
against whole-body absorbed power $\xx^H \Qmat \xx / P$.
\Cref{fig:ecbf-pareto} shows the result. The ECBF family traces the
Pareto frontier. MRT sits at maximum signal and maximum absorbed
power. The generalised-eigenvalue optimum
$\xx \propto \Qmat^{-1} \hh^*$ sits at minimum
absorbed-power-per-signal-power. The random precoder cloud and the
incoherent baseline (random-phase superposition) lie inside the
Pareto frontier as expected.

\begin{figure}[!t]
  \centering
  \includegraphics[width=\linewidth]{ecbf_pareto.pdf}
  \caption{Pareto trade-off between received signal and whole-body
  absorbed power for the 64-element URA at 28\,GHz near the
  Thelonious phantom. Panel (a) UE in front of phantom, panel (b)
  UE behind phantom. Blue curve: ECBF family parametrised by
  $\lambda \in [10^{-3}, 10^3]$. Red star: MRT. Orange diamond:
  generalised-eigenvalue optimum $\Qmat^{-1}\hh^*$. At 50\,\% of
  MRT signal power, the absorbed power drops by a factor of 4 to
  5.}
  \label{fig:ecbf-pareto}
\end{figure}

For UE position~(a), MRT gives received signal $1.47\,P$ and
absorbed power $1.93 \times 10^{-5}\,P$. For position~(b), MRT
gives $23.8\,P$ and $4.22 \times 10^{-4}\,P$. The factor of 17
difference in absorbed power between the two positions tracks the
difference in body-coupling strength under this multipath
realisation. ECBF at $\lambda$ chosen for 50\,\% of MRT signal
power reduces $P_{\mathrm{abs}}$ to $3.7 \times 10^{-6}\,P$, a
factor of 5 below MRT at the cost of 3\,dB signal loss. The
exposure-signal alignment for the two positions is
$\rho^{(a)} = 0.53$ and $\rho^{(b)} = 0.66$, both moderate. The
signal channel $\hh$ has large overlap with the leading exposure
modes of $\Qmat$ but is not aligned with
$\mathbf{v}_1$ (which would give $\rho = 1$) or with $\ker \Qmat$
(which would give $\rho = 0$). The intermediate values quantify the
geometric coupling between signal and exposure channels.

\subsection{Body-surface hot-spot maps}

\Cref{fig:hotspot} renders $\Sab(\rr)$ on the Thelonious surface
under MRT and under ECBF at 50\,\% of MRT signal power, for UE
position~(a). MRT concentrates absorption on the phantom front in
the array's main projection. ECBF spreads the absorbed power
more uniformly. Integrated $P_{\mathrm{abs}}$ drops by a factor of
4.6 at the same SNR target. The displacement of the absorbed peak
is a redistribution rather than a translation. The body's
polarisation modes that MRT excites strongly are suppressed in
favour of modes with smaller $\lambda_k$.

\begin{figure*}[!t]
  \centering
  \includegraphics[width=\linewidth]{hotspot_pair.pdf}
  \caption{Body-surface $\Sab$ on the Thelonious phantom under MRT
  (left two panels) and under ECBF at 50\,\% of MRT signal power
  (right two panels) for UE position~(a). Front and side views.
  Logarithmic colour scale from $10^{-7}$ to $10^{-3}\,$W/m$^2$.
  $P_{\mathrm{abs}}^{\mathrm{MRT}} = 1.69 \times 10^{-5}\,$W,
  $P_{\mathrm{abs}}^{\mathrm{ECBF}} = 3.72 \times 10^{-6}\,$W: a
  factor of 4.6 reduction at the same SNR target.}
  \label{fig:hotspot}
\end{figure*}

\subsection{Spectrum and effective rank}

\Cref{thm:Q-factor} predicts that $\Qmat$ has at most $M$ nonzero
eigenvalues regardless of $N$. \Cref{fig:spectrum} shows the
cumulative eigenvalue fraction $\sum_{k=1}^r \lambda_k /
\tr(\Qmat)$ as a function of rank $r$ for four multipath densities
$N/M \in \{10, 30, 60, 100\}$ at fixed $M = 64$.

\begin{figure}[!t]
  \centering
  \includegraphics[width=\linewidth]{spectrum.pdf}
  \caption{Cumulative eigenvalue fraction of $\Qmat$ as a function
  of rank for $M = 64$ and four multipath densities. Even at
  $N/M = 100$ ($N = 6400$ paths), the rank capturing 90\,\% of
  $\tr(\Qmat)$ is at most 5. Body-mode exposure information is
  low-dimensional regardless of how many paths reach the body.}
  \label{fig:spectrum}
\end{figure}

For the canonical scenario at $N/M = 40$, the effective rank that
captures 90\,\% of $\tr(\Qmat)$ is 5 at both UE positions. The
rank that captures 99\,\% is 8 at position~(a) and 9 at
position~(b). The exposure information of a 64-element array is
concentrated in roughly 5 to 10 modes regardless of multipath
density. Rank-$r$ truncation of $\Qmat$ reduces the ECBF inverse in
\cref{eq:x-opt} from an $M \times M$ problem to an $r \times r$
problem with provable error $\lambda_{r+1} P$. For the present
scenario, an $r = 10$ truncation is exact to within 0.01\,\% of
$\tr(\Qmat)$.

The combined error on $\Qmat$ for the canonical scenario is
$\|\Delta \Qmat\|_F / \|\Qmat\|_F \lesssim 3\,\%$, well within the
5\,\% envelope. The dominant contribution is \cref{prop:approx1} on
path pairs with mean incidence above $50^\circ$. Self-terms are
exact under both approximations, so the incoherent limit of
\cref{cor:incoherent} carries no approximation error.


\section{Discussion}\label{sec:discussion}

The closed-form exposure operator rests on three modelling
assumptions stated in \cref{sec:setup} and on the two cross-term
approximations stated in \cref{prop:approx1,prop:approx2}. The
combined error budget on $\xx^H \Qmat \xx$ is $\leq 5\,\%$ for any
precoder, with diagonal terms exact. The single-wave limit recovers
the polarisation-aware Fresnel surface law. The incoherent limit
recovers the standard multi-source surface dosimetry formula. Direct ground-truth
validation against full-wave coherent FDTD on a multi-element array
near a body is currently outside the reach of single-GPU
computation at 28\,GHz, where the cell count for a full body and a
64-element array crosses $10^{12}$ per configuration. Validation in
this paper rests on three independent paths: the analytical bounds
in \cref{prop:approx1,prop:approx2}, the incoherent-limit validation
inherited from geometric dosimetry, and internal consistency on the
canonical scenario above. A direct entry-by-entry comparison of the
analytical $\Qmat$ against an FDTD-calibrated $\mathbf{S}$ on a
Hochwald-2014-style scenario is open work and should become feasible
as multi-element FDTD on full bodies moves to multi-GPU clusters.

The reflection-operator complement $\Qmat_{\mathrm{ref}}$, which
describes the body's specular scattering and admits an
eigenvector-locking proposition under pseudo-Brewster, is the
subject of follow-on work. The companion structure of $\Qmat$ as a
thermal-emission operator under Kirchhoff reciprocity, and the
joint sensing-exposure duality that follows, are also deferred.

The closed-form operator reduces a $\sim 10^{12}$-cell volumetric
simulation per device-body configuration to a path-space spectral
problem that
is linear in the body mesh count and at most $M$-dimensional in the
array. Whole-body absorbed power is a Hermitian quadratic form in
the precoder. Exposure-constrained beamforming retains its
closed-form solution. Multi-user MIMO extends without new
approximations. The ray-tracing data the system already produces
for channel estimation is the same data that produces the exposure
operator, and no FDTD step enters the loop.


\section{Conclusion}\label{sec:concl}

A closed-form exposure operator $\Qmat = \Jmat^T \Mmat \Jmat$
replaces the FDTD-calibrated SAR matrix in coherent MIMO. The
operator is computed from the same propagation paths that give the
user-equipment channel, with cross-term error $\leq 5\,\%$ from two
bounded approximations and exact diagonal terms. The path-space
factorisation handles broadband, near-field, stochastic, low-rank,
and multi-user extensions as small perturbations. The
exposure-constrained beamforming problem retains the closed-form
Ying-2015 solution, with the analytical $\Qmat$ replacing the
FDTD-calibrated $\mathbf{S}$. A 64-element URA at 28\,GHz near the
Thelonious phantom shows a 4 to 5 reduction in absorbed body power
at 50\,\% of MRT signal power, and an effective rank of 5 at
90\,\% trace energy regardless of multipath density.


\bibliographystyle{IEEEtran}
\bibliography{refs}

\end{document}
